\documentclass[12pt,oneside,openany,pagenumber=footcenter]{book}
%\documentclass[a4paper,10pt]{report}

\usepackage[utf8]{inputenc}
\usepackage{lmodern}
\usepackage[T1]{fontenc}
\usepackage{amsfonts}

\usepackage{slovak}
\usepackage{fontenc}
\usepackage{graphicx}
\usepackage{graphics}
\usepackage{graphicx}
\usepackage{hyperref}
\usepackage{makeidx}

\newtheorem{definition}{Definition}[section]
\newtheorem{HLPpoznamka}{Note}[section]
\newtheorem{HLPpriklad}{Example}[section]
\newtheorem{HLPcvicenie}[HLPpriklad]{Exercise}
\newtheorem{HLPdokaz}{Proof}[section]
\newtheorem{zadanie}{Task}[section]
\newenvironment{poznamka}{\begin{HLPpoznamka}\rm}{\end{HLPpoznamka}}
\newenvironment{example}{\begin{HLPpriklad}\rm}{\end{HLPpriklad}}
\newenvironment{cvicenie}{\begin{HLPcvicenie}\rm}{\end{HLPcvicenie}}
\newenvironment{dokaz}{\begin{HLPdokaz}\rm}{\end{HLPdokaz}}
\newtheorem{veta}{Theorem}[section]
\newtheorem{lemma}[veta]{Lemma}
\newtheorem{dosledok}[veta]{Corollary}
\newtheorem{teza}[veta]{Proposition}


\pagestyle{headings}

\bibliographystyle{unsrt}

\def\indexname{Register}
% pekne pokope definujeme potrebne udaje
\def\mftitlea{Biologicky motivované výpočtové modely}
\def\mftitle{\mftitlea}
\def\mfthesistype{Dizertačná práca}
\def\mfauthor{Michal Kováč}
\def\mfadvisor{doc. RNDr. Damas Gruska, PhD.}
\def\mfplacedate{Bratislava, 2011}


\ifx\pdfoutput\undefined\relax\else\pdfinfo{ /Title (\mftitle) /Author (\mfauthor) /Creator (PDFLaTeX) } \fi

\def\eps{\varepsilon}
\def\goodgap{\hspace{\subfigcapskip}}
\makeindex

\begin{document}
\frontmatter
\thispagestyle{empty}
\begin{minipage}{0.20\textwidth}
\includegraphics[width=0.9\textwidth]{img/comenius_half.png}
\end{minipage}
\begin{minipage}{0.79\textwidth}
\begin{center}
\sc Katedra Informatiky \\
Fakulta Matematiky, Fyziky a Informatiky \\
Univerzita Komenského, Bratislava
\end{center}
\end{minipage}

\vfill
\begin{center}
\begin{minipage}{0.8\textwidth}
\hrule
\bigskip\bigskip
\centerline{\LARGE\sc\mftitlea}
\smallskip
\centerline{(\mfthesistype)}
\bigskip
\bigskip
\centerline{\large\sc\mfauthor}
\bigskip\bigskip
\hrule
\end{minipage}
\end{center}
\vfill
{\bf Vedúci:} \mfadvisor
\hfill\mfplacedate
\eject
\eject

\thispagestyle{empty}
{~}\vspace{12cm}

{~}\vspace{12cm}

\begin{minipage}{0.25\textwidth}~\end{minipage}
\begin{minipage}{0.69\textwidth}
I hereby declare that this submission is my own work and that, to the best of my knowledge and belief, it contains no material previously published or written by another person nor material which to a substantial extent has been accepted for the award of any other degree or diploma of the university or other institute of higher learning, except where due acknowledgment has been made in the text.

\bigskip\bigskip

\hfill\hbox to 6cm{\dotfill}
\end{minipage}

\chapter*{Acknowledgements}

I am very grateful to all the people I had the privilege to work with, especially I would like to thank my advisor RNDr. Damas Gruska PhD. for his constant interest, encouragement and help that I deeply appreciate.
The present work was possible due to many beneficial consultations and intensive cooperation.

I would also like to thank to the members of Laboratory of Comparative and Functional Genomics of Eukaryotic Organelles and their head, Prof. Ľubomír Tomáška for familiarization with the work of biologists.

I thank to my family, colleagues, teammates and all the friends who supported me especially last weeks before the deadlines for their patience, support and understanding\footnote{Understanding that they will propably never understand my work;)}

% Osobitná vďaka patrí vedúcemu diplomovej práce doc. RNDr. Damasovi Gruskovi PhD. za cenné rady, námety, podnetné pripomienky a všestrannú pomoc, ktorú si hlboko vážim. Len vďaka mnohým prínosným konzultáciam a intenzívnej spolupráci som bol schopný napísať toto dielo. Nesmiem zabudnúť ani na RNDr. Branislava Rovana, CSc. a spolužiakov za to, že si na dizertačnom seminári našli čas, aby si vypočuli moju prezentáciu dizertačnej práce. Ďalšie poďakovania venujem rodičom a známym, ktorí to so mnou dokázali vydržať posledné týždne pred odovzdaním.

% !TEX root = diz.tex
\chapter*{Abstract}
\begin{description} \itemsep1pt \parskip0pt \parsep0pt
  \item[Author:] \mfauthor
  \item[Title:] \mftitle
  \item[University:] Comenius University in Bratislava
  \item[Faculty:] Faculty of Mathematics, Physics and Informatics
  \item[Department:] Department of Applied Informatics
  \item[Supervisor:] \mfadvisor
\end{description}

This work discusses the research in the membrane systems, an emerging field of natural computing. Many variants of membranes systems have already been studied, most of them uses parallel rewriting and are computationally complete. Various sequential models have been proposed, however, in many cases they are weaker than their parallel variant.

We propose some other sequential models, provide universality proofs for them and suggest further research.

\begin{description}
  \item[Keywords:] Computation models inspired by biology, Membrane systems, P systems
\end{description}

\chapter*{Abstrakt}
\begin{description} \itemsep1pt \parskip0pt \parsep0pt
  \item[Autor:] \mfauthor
  \item[Názov dizertačnej práce:] \mftitle
  \item[Škola:] Univerzita Komenského v Bratislave
  \item[Fakulta:] Fakulta matematiky, fyziky a informatiky
  \item[Katedra:] Katedra aplikovanej informatiky
  \item[Vedúci dizertačnej práce:] \mfadvisor
  \item \mfplacedate
\end{description}

V tejto práci sa zaoberáme výskumom v oblasti membránových systémov. Veľa variantov už bolo preskúmaných, väčšinou pri výpočte používajú paralelizmus a sú Turingovsky úplné. Navrhlo sa aj veľa sekvenčných modelov, ale väčšina z nich má slabšiu výpočtovú silu ako ich paralelný variant.

V druhej časti predkladáme niektoré sekvenčné modely, u ktorých dokazujeme Turingovu úplnosť a navrhujeme modely, ktoré plánujeme preskúmať v rámci dizertačnej práce.

\begin{description}
  \item[Kľúčové slová:] Výpočtové modely inšpirované biológiou, Membránové systémy, P systémy
\end{description}
\tableofcontents{}
\listoffigures{}
\listoftables{}

\mainmatter
% !TEX root = diz.tex
\addcontentsline{toc}{chapter}{Introduction}

There are a lot of areas in the theoretical computer science that are motivated by other science fields. Computation models motivated by biology forms a large group of them. They include neural networks, computational models based on DNA evolutionary algorithms, which have already found their use in computer science and proved that it is worth to be inspired by biology. L-systems are specialized for describing the growth of plants, but they have also found the applications in computer graphics, especially in fractal geometry. Other emerging areas are still awaiting for their more significant uses.

One of them is the membrane computing \cite{Paun10OxfordHandbookMembraneComputing}. It is relatively young field of natural computing - in comparison: neural networks have been researched since 1943 and membrane systems since 1998 \cite{Paun98}.

Membrane systems (P systems) are distributed parallel computing devices inspired by the structure and functionality of cells. Recently, many P system variants have been developed in order to simulate the cells more realistically or just to improve the computational power.

We will start by an introduction of some natural computing areas including models inspired by biology in Chapter \ref{cha:natural_computing}. In Chapter \ref{cha:preliminaries} we recall some computer science basic notions that we will use through the work. P systems are formally presented in Chapter \ref{cha:p_systems}, with the current state of the research in their variants, overview of software simulator MeCoSym and various case studies.

In Chapter \ref{cha:on_the_edge_of_universality_of_sequential_p_systems} we will present the current state of our work, mainly from theoretic viewpoint (computational power, decidability of behavioral properties), including the published results in sections \ref{sec:inhibitors}, \ref{sec:active_membranes} and \ref{sec:notions_from_reaction_systems}.


\chapter{Preliminaries} % (fold)
\label{cha:preliminaries}

\section{Formal languages theory} % (fold)
\label{sec:formal_languages_theory}

Our study is based on the classical theory of formal languages. We will recall some definitions:

\begin{definition}
An {\bf alphabet} is a finite nonempty set of symbols. Usually it is denoted by $\Sigma$ or $V$.
\end{definition}

\begin{definition}
A {\bf string} over an alphabet is a finite sequence of symbols from alphabet.
\end{definition}

We denote by $V^*$ the set of all strings over an alphabet $V$. By $V^+$ = $V^* - \{\eps\}$ we denote the set of all nonempty strings over V.

\begin{definition}
A {\bf language} over the alphabet $V$ is any subset of $V^*$.
\end{definition}

% section formal_languages_theory (end)

\section{Register machines} % (fold)
\label{sec:register_machines}

% We will use the notion of register machine as defined in our article

\begin{definition}
  A {\bf $n$-register machine} is a tuple $M = (n,P,i,h)$, where:
  \begin{itemize}
    \item $n$ is the number of registers,
    \item $P$ is a set of labeled instructions of the form $j : (op(r),k,l)$, where $op(r)$ is an operation on register $r$ of $M$, and $j$, $k$, $l$ are labels from the set $Lab(M)$ (which numbers the instructions in a one-to-one manner),
    \item $i$ is the initial label, and
    \item $h$ is the final label.
  \end{itemize}
\end{definition}

The machine is capable of the following instructions:
\begin{itemize}
\item $(add(r),k,l)$ : Add one to the contents of register $r$ and proceed to instruction $k$ or to instruction $l$; in the deterministic variants usually considered in the literature we demand $k = l$.
\item $(sub(r),k,l)$ : If register $r$ is not empty, then subtract one from its contents and go to instruction $k$, otherwise proceed to instruction $l$.
\item $halt$ : This instruction stops the machine. This additional instruction can only be assigned to the final label $h$.
\end{itemize}

A deterministic $m$-register machine can analyze an input $(n_1,\dots,n_m)\in N_0^m$ in registers 1 to $m$, which is recognized if the register machine finally stops by the halt instruction with all its registers being empty (this last requirement is not necessary). If the machine does not halt, the analysis was not successful.

% section register_machines (end)

\section{Multisets} % (fold)
\label{sec:multisets}

\begin{definition}
A multiset over a set $X$ is a mapping $M: X\rightarrow \mathbb N$.
\end{definition}

We denote by $M(x), x\in X$ the multiplicity of $x$ in the multiset $M$.

\begin{definition}
The {\bf support} of a multiset $M$ is the set $supp(M)=\{x\in X|M(x)\geq 1\}$.
\end{definition}

It is the set of items with at least one occurence.

\begin{definition}
A multiset is {\bf empty} when it's support is empty.
\end{definition}

A multiset $M$ with finite support $X = \{x_1, x_2, \dots, x_n\}$ can be represented by the string $x_1^{M(x_1)}x_2^{M(x_2)}\dots x_n^{M(x_n)}$.

\begin{definition}
Multiset inclusion. We say that mutiset $M_1$ is included in multiset $M_2$ if $M_1(x)\leq M_2(x)\forall x \in X$. We denote it by $M_1\subseteq M_2$.
\end{definition}

\begin{definition}
The {\bf union} of two multisets $M_1\cup M_2 : X\rightarrow \mathbb N$ is defined as $(M_1\cup M_2)(x)=M_1(x)+M_2(x)$.
\end{definition}

\begin{definition}
The {\bf difference} of two multisets $M_1-M_2 : X\rightarrow \mathbb N$ is defined as $(M_1-M_2)(x)=M_1(x)-M_2(x)$.
\end{definition}

\begin{definition}
Product of multiset $M$ with natural number $n\in \mathbb N$ is $(n\cdot M)(x)=n\cdot M(x)$.	
\end{definition}

% section multisets (end)

\section{Multiset languages} % (fold)
\label{sec:multiset_languages}

The number of occurences of a given symbol $a\in V$ in the string $w\in V^*$ is denoted by $|w|_a$.

\begin{definition}
$\Psi_V(w)=(|w|_{a_1},|w|_{a_2},\dots,|w|_{a_n})$ is called a Parikh vector associated with the string $w\in V^*$, where $V=\{a_1,a_2,\dots a_n\}$.
\end{definition}

\begin{definition}
For a language $L\subseteq V^*$, $\Psi_V(L)=\{\Psi_V(w)|w\in L\}$ is the Parikh mapping associated with $V$.
\end{definition}

\begin{example}
Consider an alphabet $V=\{a,b\}$ and a language $L=\{a, ab, ba\}$.
$\Psi_V(L)=\{(1,0), (1,1)\}$. Notice that Parikh image of $L$ has only 2 element while $L$ has 3 elements.
\end{example}

\begin{definition}
If FL is a family of languages, by PsFL is denoted the family of Parikh images of languages in FL.
\end{definition}


% section multiset_languages (end)

Parikh images

% chapter preliminaries (end)

\chapter{Membrane computing}

Membranes intro

Natural computing is a recent field of research wchi tries to imitate nature in the way it "computes", learning new computing models and computing paradigms experimented for billions of years by nature.

Neural networks, genetic algorithms and DNA computing are already well established research fields.

However, nature computes not only at the neural or genetic level, but also at the cellular level. In general, any non-trivial biological system has a hierarchical structure where objects and information flows between regions, what can be interpreted as a computation process.

% The notion of membrane

The regions are typically delimited by various types of membranes at different levels from cell membranes, through skin membrane to virtual membranes which delimits different parts of an ecosystem.
This hierarchical system can be seen in other field such as distributed computing, where again well delimited computing units coexist and are hierarchically arranged in complex systems from single processors to the internet.

Membranes keep together certain chemicals or information and selectively determines which of them may pass through.

% The notion of membrane structure

From these observations, Paun \cite{Paun98} introduces the notion of a membrane structure as a mathematical representation of hierarchical architectures composed of membranes. It is usually represented as a Venn diagram with all the considered sets being subsets of a unique set and not allowed to be intersected. Every two sets are either one the subset of the other, or disjoint.






P systems

\chapter{P systems}

In previous chapter we introduced the notions of membrane and membrane structure.

% Place objects in the regions.

The next step is to place certain objects in the regions delimited by the membranes. The objects are identified by their names, mathematically symbols from a given alphabet.

% Multisets of objects.

Several copies of the same object can appear in a region, so we will work with multisets of objects.

% Evolution rules

In order to obtain a computing device, we will allow the objects to evolve according to evolution rules. Any object, alone or together with another objects, can be transformed in other objects, can pass through a membrane, and can dissolve the membrane in which it is placed.

% Parallelism

All objects evolve at the same time, in parallel manner across all membranes.

% Priorities

The evolution rules are hierarchizes by a priority relation, which is a partial order.

% P system

These aspects all together forms a P system as introduced in \cite{Paun98}.

In section~\ref{sec:p-systems-definitions} we will provide formal definition of a P system.

\section{Definitions}
\label{sec:p-systems-definitions}

% Definition taken from my article

{\bf P system} is a tuple $(V, \mu, w_1, w_2,\dots , w_m, R_1, R_2, \dots , R_m)$, where:
\begin{itemize}
  \item $V$ is the alphabet of symbols,
  \item $\mu$ is a membrane structure consisting of $m$ membranes labeled with numbers $1,2,\dots,m$,
  \item $w_1,w_2,\dots w_m$ are multisets of symbols present in the regions $1,2,\dots,m$ of the membrane structure,
  \item $R_1,R_2,\dots R_m$ are finite sets of the rewriting rules associated with the regions $1,2,\dots,m$ of the membrane structure.
\end{itemize}


Each rewriting rule may specify for each symbol on the right side, whether it stays in the current region, moves through the membrane to the parent region or through membrane to one of the child regions. An example of such rule is the following: $abb\rightarrow (a,here)(b,in)(c,out)(c,here)$.

% Configuration

A {\bf configuration} of a P system is represented by its membrane structure and the multisets of objects in the regions.

% Step

A {\bf computation step} of P system is a relation $\Rightarrow$ on the set of configurations such that $C_1 \Rightarrow C_2$ iff:

For every region in $C_1$ (suppose it contains a multiset of objects $w$) the corresponding multiset in $C_2$ is the result of applying a multiset of maximal simultaneously applicable multiset rewriting rules in $R^{msap}_w$ to $w$.

In other words, a maximal multiset of rules is applied in each region.

For example, let's have two regions with multisets $aa$ and $b$. In the first region there is a rule $a\rightarrow b$ and in the second membrane there is a rule $b\rightarrow aa$. The only possible result of a computation step is $bb$, $aa$. The first rule was applied twice and the second rule once. No more object could be consumed by rewriting rules.

% Computation

{\bf Computation} of a P system consists of a sequence of steps. The step $S_i$ is appied to result of previous step $S_{i-1}$. So when $S_i = (C_j,C_{j+1})$, $S_{i-1} = (C_{j-1},C_j)$.

% Result of a computation

There are two possible ways of assigning a result of a computation:

\begin{enumerate}
    \item By considering the multiplicity of objects present in a designated membrane in a halting configuration. In this case we obtain a vector of natural numbers. We can also represent this vector as a multiset of objects or as Parikh image of a language.
    \item By concatenating the symbols which leave the system, in the order they are sent out of the skin membrane (if several symbols are expelled at the same time, then any ordering of them is accepted). In this case we generate a language.
\end{enumerate}

The result of a computation is clearly only one multiset or a string, but for one initial configuration there can be multiple possible computations. It follows from the fact that there exist more than one maximal multiset of rules that can be applied in each step.

TODO:

How is P system defined. Mostly taken from my article. Multiset rewriting. Accepting vs generating model, active or passive membranes, quality vs quantity aspects.

\section{P system variants}
We are interested in computation power of various variants of P systems. Especially those that are universal (Turing completeness).

\subsection{Parallelism options}
Maximal parallelism, minimal parallelism, n-parallelism, sequential models.

\subsection{Contextivity rules}
Context rules vs cooperational rules, catalytic rules, symmetric cooperational rules, catalytic rules, promoters, inhibitors, context-free rules.

\subsection{Priority rules}
\subsection{Energy of membranes}
\subsection{Calculi of Looping Sequences}

\section{Case studies}
Vultures in Pyrenees, Scavangers of Pyrenees.

\chapter*{Conclusions}
\addcontentsline{toc}{chapter}{Conclusions}
We have studied several variants of sequential P systems in order to obtain universality without using maximal parallelism. A variant with rewriting rules that can use inhibitors was shown to be universal in both generating and accepting case. The generating model is able to simulate maximal parallel P system and the accepting model can simulate a register machine.

In addition, we have defined a new notion of vacuum, which is immediately created in the region that becomes empty. A new P system variant was introduced, which allowed the vacuum to be used on the left side of P system rewriting rules. The accepting case of this variant was shown to be universal by direct implementation of a register machine.

A lot of other variants deserve to be combined with the vacuum, so we suggest to research them more
thoroughly(non-cooperative rules, rules with priorities, decaying objects, deterministic steps, . . . ).
\backmatter

% \begin{thebibliography}{1}
\bibliography{diz}
% \end{thebibliography}

\appendix

\chapter{Štatistiky}
Tu budú štatistiky.

\printindex

\end{document}
